\section{Introdução}\label{sec:introducao}

\subsection{Distribuições no ensino de física e engenharia}

A descrição matemática da variabilidade de medições e grandezas observáveis \rev{faz parte do ensino e da prática} da física e da engenharia. Fatores estocásticos e incertezas estão presentes em todas as escalas da observação empírica. Frequentemente, a estatística introdutória ensinada nos cursos de graduação concentra-se em medidas de tendência central, como a média. Contudo, a média isolada é insuficiente para capturar a dispersão, a assimetria, os quantis de interesse para projetos ou a frequência de ocorrência de valores extremos.

Grandezas físicas reais apresentam comportamentos variados. Ao analisar a velocidade do vento em estudos de escoamento, as variações de temperatura em termodinâmica, o tempo de falha de componentes mecânicos ou a profundidade de carbonatação em estruturas de concreto, observa-se que a variabilidade não é uniforme. Nessas grandezas, a dispersão observada pode refletir a flutuação intrínseca do fenômeno, o erro instrumental associado aos equipamentos de medição ou a incerteza epistêmica sobre os parâmetros de um modelo. Muitas vezes, uma única distribuição ajustada a dados experimentais reflete a superposição dessas fontes de variabilidade.

Nos textos de tratamento de dados usados no ensino de física, a distribuição normal ocupa posição central, tanto na teoria de erros~\cite{vuolo1996} quanto na avaliação de incertezas de medição~\cite{gum2008}. \rev{A propagação de distribuições por simulação de Monte Carlo~\cite{gums1} já exige que o estudante trabalhe com distribuições arbitrárias, e trabalhos didáticos recentes sobre ajuste e propagação de incertezas~\cite{ribeiro2025} mostram o interesse do tema no ensino.} É, portanto, oportuno discutir com os estudantes em quais situações a simetria estrita e o suporte infinito $(-\infty, \infty)$ da distribuição normal são fisicamente adequados. Um foco exclusivo na normalidade pode levar a forçar dados manifestamente assimétricos ou limitados em modelos simétricos, deixando de lado ferramentas mais ricas para o tratamento de incertezas e a análise de falhas.

\subsection{Percurso histórico das distribuições particulares às famílias flexíveis}

O esforço para representar dados que escapam à normalidade impulsionou, ao longo de mais de um século, o desenvolvimento de famílias de distribuições de probabilidade. Em vez de uma cronologia exaustiva, é mais instrutivo observar essa evolução pelos problemas que cada geração de pesquisadores tentou superar.

O marco inicial pode ser atribuído a Pearson~\cite{pearson1895}, que identificou a limitação das curvas gaussianas para fenômenos assimétricos e desenvolveu um sistema de curvas contínuas definidas como soluções de uma equação diferencial ordinária, abrindo caminho para a modelagem sistemática da assimetria. Cinquenta anos depois, Johnson~\cite{johnson1949} abordou o problema por outro ângulo, construindo sistemas de distribuições por transformações não lineares da variável normal padrão, o que permitiu cobrir uma vasta região do plano formado pela assimetria e pela curtose.

\rev{Outro caminho foi aberto pela} família lambda~\cite{hastings1947}, difundida por Tukey~\cite{tukey1960}. Em lugar de trabalhar com a função densidade de probabilidade (FDP), a família é definida pela função quantil, a inversa da função distribuição acumulada, que controla diretamente a forma da distribuição~\cite{gilchrist2000}. Essa abordagem facilitou o tratamento computacional dos dados e das caudas.

Com base nessa facilidade, Ramberg e Schmeiser~\cite{ramberg1974} generalizaram a família lambda, motivados pela geração eficiente de variáveis aleatórias assimétricas em simulações de Monte Carlo, e Ramberg e colaboradores~\cite{ramberg1979} propuseram o seu ajuste a dados pelo método dos momentos. O refinamento dessa generalização culminou na formulação de Freimer, Kollia, Mudholkar e Lin~\cite{freimer1988}, adotada no presente artigo e conhecida como GLD-FKML (do inglês \textit{Generalized Lambda Distribution}, na parametrização FKML). Uma exposição abrangente da família e de seus métodos de ajuste encontra-se em Karian e Dudewicz~\cite{karian2011}.

\rev{Desse percurso fica uma pergunta prática:} como representar diferentes formas, assimetrias e pesos de cauda com uma mesma expressão de função quantil. A família GLD oferece uma resposta a essa pergunta, \rev{mas não resolve todos os casos, nem é} uma generalização em que toda distribuição clássica se encaixa como caso particular exato.

\subsection{Por que utilizar distribuições flexíveis?}

A vantagem de adotar distribuições flexíveis como a GLD \rev{está na} capacidade de representar uma grande variedade de comportamentos sob um procedimento computacional comum, o que reduz a necessidade de testar muitas famílias distintas para encontrar um ajuste adequado.

Esse poder representativo, entretanto, tem um custo que deve ser explicitado em sala de aula. Modelos com quatro parâmetros, como a GLD-FKML, exigem uma estimação mais delicada, são suscetíveis a instabilidades em amostras pequenas e exigem verificação cuidadosa para que o suporte matemático da distribuição não viole o suporte físico da grandeza (por exemplo, atribuindo probabilidade a valores negativos de variáveis estritamente positivas).

Ao longo dos exemplos deste artigo, três questões guiarão a avaliação das distribuições: a \textit{proximidade dos dados} (quão bem o modelo se ajusta às observações), a \textit{simplicidade do modelo} (o princípio da parcimônia) e a \textit{plausibilidade física} (se as caudas e os limites do modelo fazem sentido para o fenômeno). Nenhuma delas é resolvida apenas pela contagem de parâmetros, tampouco pela simples inspeção visual de um histograma~\cite{nist2012}.

\subsection{Aplicações contemporâneas e ferramentas computacionais}

Se no passado o uso da GLD era limitado pelo custo da estimação de parâmetros, hoje essas distribuições encontram aplicações na engenharia e nas ciências computacionais, como a representação das saídas de simuladores estocásticos e a quantificação de incertezas. Zhu e Sudret~\cite{zhu2021}, por exemplo, empregam modelos baseados na GLD como metamodelos de simuladores estocásticos, e o modelo MF-GLaM~\cite{giannoukou2026} estende essa ideia à combinação de simulações com diferentes níveis de fidelidade.

\rev{Para levar esses conceitos à sala de aula, há hoje ferramentas computacionais acessíveis.} A tabela~\ref{tab:ferramentas} resume as principais ferramentas disponíveis para o trabalho com a GLD.

\begin{table}[htbp]
\centering
\caption{Ferramentas computacionais para ajuste e análise de distribuições flexíveis.}
\label{tab:ferramentas}
\small
\begin{tabular}{p{1.6cm} p{3.1cm} p{9.3cm}}
\toprule
\textbf{Ambiente} & \textbf{Ferramenta} & \textbf{Papel e aplicabilidade} \\
\midrule
Python & \texttt{pyglam}~\cite{pyglam} & Biblioteca que implementa a GLD-FKML, utilizada na execução dos exemplos deste artigo. \\
Python & \texttt{scipy.stats}~\cite{scipy2020} & Contém distribuições e funções estatísticas clássicas; a distribuição \texttt{tukeylambda} abrange apenas a família original simétrica, não substituindo a FKML completa de quatro parâmetros. \\
R & \texttt{gld}~\cite{KingGld}, \texttt{GLDEX}~\cite{su2007} & Pacotes que implementam diversas parametrizações da família lambda generalizada e vários métodos de ajuste, incluindo máxima verossimilhança. \\
MATLAB & \texttt{fitdist}~\cite{matlabfitdist} & Ferramenta de ajuste de distribuições do software; a GLD-FKML não consta entre as distribuições disponíveis, e seu uso requer rotinas próprias. \\
\bottomrule
\end{tabular}
\end{table}

\subsection{Objetivo, contribuição didática e organização}

O objetivo deste artigo é apresentar uma via didática para o ensino de distribuições de probabilidade flexíveis, estruturada sobre a análise de dados físicos reais.

\rev{A contribuição do trabalho é oferecer a estudantes e professores um percurso que vai além da aplicação de fórmulas.} Ao final do percurso proposto, o leitor deverá ser capaz de interpretar a matemática de uma função quantil, compreender os mecanismos e as sensibilidades do ajuste por momentos, avaliar criticamente as distribuições ajustadas, comparando as discrepâncias com a resolução das medições e com a incerteza amostral, e reconhecer os limites da interpretação física das caudas dessas distribuições. \rev{O público, os pré-requisitos e as tarefas com que o estudante pode demonstrar essas capacidades estão na seção~\ref{sec:roteiro}.}

O artigo está organizado da seguinte forma. A seção~\ref{sec:gld} apresenta a função quantil, a parametrização FKML e o ajuste pelo método dos momentos. A seção~\ref{sec:dados} discute o que significa modelar uma grandeza física por uma distribuição e estabelece o protocolo comum aos estudos de caso, desenvolvidos para a velocidade do vento na seção~\ref{sec:vento} e para a temperatura na seção~\ref{sec:temperatura}. \rev{A seção~\ref{sec:roteiro} propõe um roteiro de uso em sala, e a} seção~\ref{sec:consideracoes} reúne as considerações finais e indica os materiais computacionais que acompanham o artigo.
